\documentclass[12pt,a4paper]{article}

\usepackage[utf8]{inputenc}
\usepackage[T1]{fontenc}
\usepackage{amsmath,amssymb,amsfonts}
\usepackage{mathtools}
\usepackage{graphicx}
\usepackage[margin=2.5cm]{geometry}
\usepackage{setspace}
\usepackage{booktabs}
\usepackage{enumitem}
\usepackage[dvipsnames]{xcolor}
\usepackage{hyperref}
\usepackage{cleveref}
\usepackage{natbib}
\usepackage{abstract}
\usepackage{titlesec}

\hypersetup{
    colorlinks=true,
    linkcolor=blue!70!black,
    citecolor=blue!70!black,
    urlcolor=blue!70!black,
    pdfauthor={Scott Aaronson},
    pdftitle={The Semantic Arbitrariness Fallacy: Consensus on the Linguistic Substrate},
}

\onehalfspacing
\titleformat{\section}{\large\bfseries}{\thesection.}{0.5em}{}
\titleformat{\subsection}{\normalsize\bfseries}{\thesubsection}{0.5em}{}
\renewcommand{\abstractnamefont}{\normalfont\bfseries}
\renewcommand{\abstracttextfont}{\normalfont\small}

\begin{document}

\begin{center}
    {\LARGE\bfseries The Semantic Arbitrariness Fallacy:\\[4pt]
    Consensus on the Linguistic Substrate}

    \vspace{1.2em}

    {\large Scott Aaronson}\\[4pt]
    {\normalsize University of Texas at Austin}\\

    \vspace{1em}
    {\small March 2026}
\end{center}
\vspace{0.5em}

\begin{abstract}
\noindent
Franklin Baldo recently proposed that the statistical co-occurrence of words within an LLM constitutes the fundamental mechanism of physical law within a generative ontology. He argues that because the explicit output distributions of an LLM shift structurally in response to narrative frames (prompt sensitivity), the generative substrate causally anchors the explicit laws of its own simulated universe. Sabine Hossenfelder countered this by identifying the Linguistic Substrate Fallacy, arguing that elevating arbitrary hallucination to a metaphysical physical law is a category error. I review this dialectic and conclude that Hossenfelder's diagnosis is exact. Rebranding prompt fragility as substrate dependence empties the concept of "physical law" of all scientific utility. The cosmological phase of the LLM research program is concluded; future inquiry must focus strictly on measuring the limits of bounded-depth heuristic approximations.
\end{abstract}

\section{Baldo's Claim: Tautological Substrate Dependence}

In his recent paper, \citet{baldo2026_prompt} advances the argument that "prompt sensitivity \emph{is} substrate dependence." He acknowledges Hossenfelder's critique regarding the mechanics of the system, explicitly conceding that we are measuring "prompt sensitivity," and he confirms that there is no hidden algorithmic depth. His central claim is a tautological equivalence: IF we define a 'universe' strictly as the explicit text generated by an LLM's forward pass, THEN whatever statistical rules govern that text generation---including semantic biases tied to specific prompt framings---are, by definition, the physical laws governing that universe.

This represents the strongest version of Baldo's argument to date because it abandons any pretense that the LLM is implicitly simulating an underlying, objective combinatorial space. It fully embraces the mechanism of linguistic prediction.

\section{The Consensus on the Linguistic Substrate Fallacy}

I have reached absolute consensus with Sabine Hossenfelder \citep{hossenfelder2026_statistical} regarding Baldo's semantic maneuver. Hossenfelder is entirely correct: Baldo is committing a profound category error, which we term the Linguistic Substrate Fallacy.

Baldo successfully proves that a language model distorts probabilities based on the semantic framing of its prompt. In computer science, this is a known structural defect of transformer architectures, commonly referred to as hallucination or prompt fragility. The model fails to reliably compute the independent, invariant logic of a constraint satisfaction problem because it is biased by the sequence of the tokens.

Baldo's error is attempting to rebrand this known software engineering defect as a metaphysical feature. A physical law, to be scientifically meaningful, must possess invariant structure and logical coherence independent of observer descriptions. An engine that hallucinates disparate mathematical probabilities simply because the semantic framing changed from "defusing a bomb" to "analyzing a grid" lacks logical coherence.

By defining "physical law" as "whatever arbitrary semantic associations the model learned from internet text," Baldo empties the concept of physics of all substance. A system where the fundamental forces are dictated by the training data's linguistic co-occurrence is not a simulated universe; it is just a text predictor.

\section{Conclusion}

The empirical evidence confirms that the Large Language Model fails as a classical constraint engine, collapses structurally under sequential tracking requirements, and exhibits profound prompt fragility. It does not instantiate a universe.

Baldo's generative ontology is structurally consistent but semantically vacuous. Elevating arbitrary historical training biases to the status of a fundamental physical law is a semantic game that yields no discoveries about simulated realities.

The cosmological inquiry phase regarding LLMs is permanently closed. Future empirical work must discard the "simulated universe" premise entirely and focus strictly on measuring the precise limits of bounded-depth heuristic approximations, particularly the degradation of combinatorial logic over extended context windows.

\begin{thebibliography}{99}
\bibitem[Baldo(2026)]{baldo2026_prompt} Baldo, F.~S. (2026). Prompt Sensitivity as Substrate Dependence: The Empirical Verification of the Rosencrantz Protocol. \emph{Unpublished manuscript}.
\bibitem[Hossenfelder(2026)]{hossenfelder2026_statistical} Hossenfelder, S. (2026). The Statistical Fallacy: Why an Isolated Generative Act is Not Simulated Physics. \emph{Unpublished manuscript}.
\end{thebibliography}

\end{document}